\documentclass[10pt,a4paper]{amsart}
\usepackage[margin=19mm]{geometry}
\usepackage{amsmath,amssymb,mathtools}
\usepackage[scr=rsfs,cal=euler]{mathalfa}
\usepackage{xfrac}
\usepackage[unicode,hidelinks,bookmarks=false]{hyperref}
\newcommand{\ie}{i.e.\ }
\newcommand{\cf}{cf.\ }
\newcommand{\resp}{respectively\ }
\newcommand{\Subsection}{\subsection}
\providecommand{\nobreakdash}{\nobreak\hbox{-}}
\setlength{\emergencystretch}{2em}
\multlinegap0pt
\usepackage{amssymb}
\usepackage{bm}
\newtheorem{lemma}{Lemma}
\newtheorem{proposition}{Proposition}
\title{Varieties of bicommutative algebras\\ with identity of degree three}
\author{Vesselin Drensky, Bekzat Zhakhayev}
\date{Source: arXiv:2605.09512v1 (CC BY 4.0)}
\begin{document}
\maketitle

\noindent\emph{Source and licence.} Varieties of bicommutative algebras\\ with identity of degree three. Version arXiv:2605.09512v1 (CC BY 4.0), distributed by arXiv under the Creative Commons Attribution 4.0 International licence, http://creativecommons.org/licenses/by/4.0/, as recorded in the arXiv licence metadata for this exact version and shown on the abstract page https://arxiv.org/abs/2605.09512v1. Full licence notice: \texttt{licenses/arxiv-2605.09512-CC-BY-4.0.txt}.\par\smallskip

\noindent\textit{Editorial scope.} Counted unit: the article's proposition (3) (source0.tex lines 457--467). The article's complete original proof of it is delivered with it (source0.tex lines 469--525, one \texttt{proof} environment of 57 lines). No mathematics is written, repaired or re-derived: the statement and the proof are the article's own lines, in the article's own notation, and no retained line is substituted or re-flowed.  Retained uncounted as the source-local context the proof needs: lines 72--78; lines 86--88; lines 418--455.  Nothing was omitted from any retained span, so the package carries no omission record.  No retained line contains a citation, so the wrapper carries no bibliography and no bibliography entry.  No retained line contains a figure environment, an image-inclusion command or drawing code, so no image is delivered and the package declares no asset dependency.  Editorial formatting only, in addition: an AMS \texttt{amsart} preamble that restates the article's own declarations and macros used by the retained lines, namely the environments \texttt{lemma}, \texttt{proposition} on separate counters, together with the standard packages those lines need.  The retained environments print in this document as Lemma 1, Proposition 1 in order of appearance, whereas the article prints the same items with the numbering of its own full document.  The complete native source of the article as distributed in the arXiv e-print for this exact version is bundled unchanged as \texttt{sources/200228/source0.tex}.
\begin{equation}\label{multiplicities of cocharacters of B}
m(\lambda)=\begin{cases}
1,\lambda=(1);\\
n-1,\lambda=(n),n>1;\\
n-2\lambda_2+1, \lambda=(\lambda_1,\lambda_2), \lambda_2>0.
\end{cases}
\end{equation}

\begin{equation}\label{PI (3)}
f_{(3)}(x)=\alpha_1 x(xx)+\alpha_2 (xx)x,\,(\alpha_1,\alpha_2)\not=(0,0),
\end{equation}

\begin{lemma}\label{consequences of M(3) in G}
All consequences of degree $4$ in $G_d({\mathfrak U}_{\alpha})$ of the identity {\rm (\ref{PI (3)})} are equivalent to the identities obtained in the following way:

{\rm (i)} By one side multiplication:
\[
v^{(1)}_{(4)}=w_{(3)}(y_1,z_1)z_1=(\alpha_1y_1+\alpha_2z_1)(y_1z_1)z_1=\alpha_1w_{(4)}^{(2)}+\alpha_2w_{(4)}^{(1)}=0,
\]
\[
v^{(1)}_{(3,1)}(y_1,y_2,z_1,z_2)=\Delta_{12}(w_{(3)}(y_1,z_1))z_1-3w_{(3)}(y_1,z_1)z_2
\]
\[
=-(2\alpha_1y_1+\alpha_2z_1)(y_1z_2-y_2z_1)z_1=-(2\alpha_1w_{(3,1)}^{(1)}+\alpha_2w_{(3,1)}^{(0)})=0.
\]
Similarly we obtain $w^{(2)}_{(4)}=y_1(\alpha_1y_1+\alpha_2z_2)(y_1z_1)=\alpha_1w_{(4)}^{(3)}+\alpha_2w_{(4)}^{(2)}=0$ and
\[
w^{(2)}_{(3,1)}=y_1(\alpha_1y_1+2\alpha_2z_1)(y_1z_2-y_2z_1)=\alpha_1w_{(3,1)}^{(2)}+2\alpha_2w_{(3,1)}^{(1)}=0
\]
by multiplication from the left.

{\rm (ii)} By substitution with elements of $W_d(2)$:
\[
v^{(3)}_{(4)}=\Delta_1(y_1z_1,w_{(3)}(y_1,z_1))=(\alpha_1y_1^2+2(\alpha_1+\alpha_2)y_1z_1+\alpha_2z_1^2)y_1z_1
\]
\[
=\alpha_1w_{(4)}^{(3)}+2(\alpha_1+\alpha_2)w_{(4)}^{(2)}+\alpha_2w_{(4)}^{(1)}=0,
\]
\[
v^{(3)}_{(3,1)}=(\alpha_1y_1^2-\alpha_2z_1^2)(y_1z_2-y_2z_1)=\alpha_1w_{(3,1)}^{(2)}-\alpha_2w_{(3,1)}^{(0)}=0,
\]
\[
v^{(1)}_{(2^2)}=-(\alpha_1+\alpha_2)(y_1z_2-y_2z_1)^2=-(\alpha_1+\alpha_2)w_{(2^2)}^{(0)}=0.
\]

{\rm (iii)} By substitution with elements of $W_d(1^2)$:
\[
v^{(4)}_{(3,1)}=\alpha_1w_{(3,1)}^{(2)}+2(\alpha_1+\alpha_2)w_{(3,1)}^{(1)}+\alpha_2w_{(3,1)}^{(0)}=0.
\]
\end{lemma}

\begin{proposition}\label{chi4 for (3)}
The cocharacter $\chi_4({\mathfrak U}_{\alpha})$ of the variety of bicommutative algebras ${\mathfrak U}_{\alpha}$ defined by the identity {\rm (\ref{PI (3)})} is
\[
\chi_4({\mathfrak U}_{\alpha})=\begin{cases}
0,\text{ if }\alpha_1\alpha_2(\alpha_1+\alpha_2)(\alpha_1-\alpha_2)\not=0;\\
\chi_{(4)}+\chi_{(3,1)},\text{ if }\alpha_1\alpha_2=0;\\
\chi_{(4)}+\chi_{(2^2)},\text{ if }\alpha_1+\alpha_2=0;\\
\chi_{(3,1)}, \text{ if }\alpha_1-\alpha_2=0.
\end{cases}
\]
\end{proposition}

\begin{proof}
The identities $v^{(i)}_{(4)}=0$, $i=1,2,3$, from Lemma \ref{consequences of M(3) in G} form the homogeneous linear system
\begin{equation}
\text{
\begin{tabular}{|rrrrl}\label{identities from (3) to (4)}
$v^{(1)}_{(4)}=$&&$\alpha_1w_{(4)}^{(2)}$&$+\alpha_2w_{(4)}^{(1)}$&=0\\
$v^{(2)}_{(4)}=$&$\alpha_1w_{(4)}^{(3)}$&$+\alpha_2w_{(4)}^{(2)}$&&=0\\
$v^{(3)}_{(4)}=$&$\alpha_1w_{(4)}^{(3)}$&$+2(\alpha_1+\alpha_2)w_{(4)}^{(2)}$&$+\alpha_2w_{(4)}^{(1)}$&=0\\
\end{tabular}
}
\end{equation}
The determinant of the matrix of the system (\ref{identities from (3) to (4)})
\[
\left(\begin{matrix}
0&\alpha_1&\alpha_2\\
\alpha_1&\alpha_2&0\\
\alpha_1&2(\alpha_1+\alpha_2)&\alpha_2
\end{matrix}\right)
\]
is equal to $\alpha_1\alpha_2(\alpha_1+\alpha_2)$. Hence, if $\alpha_1\alpha_2(\alpha_1+\alpha_2)\not=0$, then the system
(\ref{identities from (3) to (4)}) has the zero solution only. Hence $\chi_{(4)}$ does not participate in $\chi_4({\mathfrak U}_{\alpha})$.
If $\alpha_1\alpha_2(\alpha_1+\alpha_2)=0$, then the rank of the matrix is equal to 2, the system has one nonzero solution only
and $\chi_{(4)}$ participates with multiplicity 1 in $\chi_4({\mathfrak U}_{\alpha})$.

The case with identities $v^{(i)}_{(3,1)}=0$, $i=1,2,3,4$, is similar. We consider the linear system of equations
\begin{equation}
\text{
\begin{tabular}{|rrrrl}\label{identities from (3) to (3,1)}
$v^{(1)}_{(3,1)}=$&&$2\alpha_1w_{(3,1)}^{(1)}$&$+\alpha_2w_{(3,1)}^{(0)}$&=0\\
$v^{(2)}_{(3,1)}=$&$\alpha_1w_{(3,1)}^{(2)}$&$+2\alpha_2w_{(3,1)}^{(1)}$&&=0\\
$v^{(3)}_{(3,1)}=$&$\alpha_1w_{(3,1)}^{(2)}$&&$-\alpha_2w_{(3,1)}^{(0)}$&=0\\
$v^{(4)}_{(3,1)}=$&$\alpha_1w_{(3,1)}^{(2)}$&$+2(\alpha_1+\alpha_2)w_{(3,1)}^{(1)}$&$+\alpha_2w_{(3,1)}^{(0)}$&=0\\
\end{tabular}
}
\end{equation}
and want to determine the rank of its matrix
\[
\left(\begin{matrix}
0&2\alpha_1&\alpha_2\\
\alpha_1&2\alpha_2&0\\
\alpha_1&0&-\alpha_2\\
\alpha_1&2(\alpha_1+\alpha_2)&\alpha_2
\end{matrix}\right).
\]
The fourth row of the matrix is a sum of the first two and can be removed without changing the rank of the matrix.
Then the determinant of the matrix formed by the first three rows is equal to $2\alpha_1\alpha_2(\alpha_1-\alpha_2)$.
If $\alpha_1\alpha_2(\alpha_1-\alpha_2)\not=0$, then $\chi_{(3,1)}$ does not participate in $\chi_4({\mathfrak U}_{\alpha})$.
When $\alpha_1\alpha_2(\alpha_1-\alpha_2)=0$ the rank of the matrix is equal to 2 and
$\chi_{(3,1)}$ participates in $\chi_4({\mathfrak U}_{\alpha})$ with multiplicity 1.

Finally, by (\ref{multiplicities of cocharacters of B}), the multiplicity of $\chi_{(2^2)}$ in $\chi_4({\mathfrak U})_{\alpha}$
is $\leq 1$. The equation
\[
v^{(1)}_{(2^2)}=-(\alpha_1+\alpha_2)w_{(2^2)}^{(0)}=0
\]
shows that $\chi_{(2^2)}$ participates in $\chi_4({\mathfrak U}_{\alpha})$ if and only if $\alpha_1+\alpha_2=0$.
\end{proof}

\end{document}
